\documentclass[a4paper,11pt]{article}

% === Document Encoding and Fonts ===
\usepackage[utf8]{inputenc}
\usepackage[T1]{fontenc}
\usepackage{microtype}

% === Page Layout ===
\usepackage[top = 1.2in, bottom = 1.2in, left=1in, right=1in]{geometry}

% === Core Mathematics Packages ===
\usepackage{mathtools}
\usepackage{amsthm}
\usepackage{amssymb}
\usepackage{amsfonts}
\usepackage{mathrsfs}
\usepackage{nicefrac}
\usepackage{euscript}

% === Diagrams and Graphics ===
\usepackage{tikz-cd}
\usepackage[unicode=true,
            colorlinks=true,
            linkcolor=blue,
            urlcolor=blue,
            citecolor=red]{hyperref}

% === Theorem Environments ===
\theoremstyle{plain}
\newtheorem{thm}{Theorem}[section]
\newtheorem{prop}[thm]{Proposition}
\newtheorem{lem}[thm]{Lemma}
\newtheorem{cor}[thm]{Corollary}

\theoremstyle{definition}
\newtheorem{defn}[thm]{Definition}
\newtheorem{exmp}[thm]{Example}
\newtheorem{constr}[thm]{Construction}

\theoremstyle{remark}
\newtheorem{rem}[thm]{Remark}

% === Custom Macros ===
\newcommand{\bbR}{\mathbb{R}}
\newcommand{\bbN}{\mathbb{N}}
\newcommand{\bbZ}{\mathbb{Z}}
\newcommand{\bbS}{\mathbb{S}}
\newcommand{\flowM}{\mathcal{M}}
\newcommand{\fC}{\mathscr{C}}
\newcommand{\sD}{\mathscr{D}}
\newcommand{\op}{\{1 \to 0\}}
\newcommand{\Image}{\operatorname{Im}}
\newcommand{\cofib}{\operatorname{cof}}
\newcommand{\colim}{\operatorname{colim}}
\newcommand{\Sp}{\mathbf{Sp}}
\newcommand{\Top}{\mathbf{Top}}
\newcommand{\Ch}{\operatorname{Ch}}
\newcommand{\Fil}{\operatorname{Fil}}
\newcommand{\gr}{\operatorname{gr}}
\newcommand{\ob}{\operatorname{ob}}
\newcommand{\ind}{\operatorname{ind}}
\newcommand{\fr}{\operatorname{fr}}
\newcommand{\norm}[1]{\left\lvert#1\right\rvert}

\date{\today}
\title{Exact Couples and Differentials of Flow Categories:\\A Geometric Perspective on Amanda Hirschi's Proposition 2.20}
\author{Arye Deutsch}

\begin{document}

\maketitle

\begin{abstract}
This note provides a self-contained, mathematically rigorous exposition of the differentials in the spectral sequence associated with a filtered flow category. Starting with the abstract algebraic machinery of exact couples, we derive the general formula for the $r$-th page differential $d_r$ and prove its well-definedness. We then apply these algebraic definitions to the geometric setting of a filtered flow category $\flowM$, as described in Proposition 2.20 of Amanda Hirschi's Lecture Notes. We show that the first-page differential $d_1$ recovers the standard Morse differential, that the survival of a class to page $r+1$ is equivalent to extending it to an $r$-truncated right flow module, and that the higher differentials $d_{r+1}$ are given geometrically by the Toda manifolds obtained by gluing these modules along the flow manifolds of $\flowM$.
\end{abstract}

\tableofcontents
\newpage

\section{Introduction}

In Floer homotopy theory, stable homotopy types of flow categories are represented by spectra whose homotopy groups can be calculated via spectral sequences. A natural way to obtain such a spectral sequence is to filter the flow category $\flowM$ by the degree of its objects:
\[
\dots \subset F^{i+1} \flowM \subset F^i \flowM \subset F^{i-1} \flowM \subset \dots
\]
where $\ob(F^i \flowM) = \{x \in \ob(\flowM) \mid |x| > i\}$. 

As established in Proposition 2.20 of Amanda Hirschi's summer school lectures, the filtration on the category induces a sequence of exact triangles on the corresponding bordism groups:
\[
\dots \to \Omega_k(F^i \flowM) \xrightarrow{\text{ext}} \Omega_k(F^{i-1} \flowM) \xrightarrow{\text{rest}} \Omega_k(F^{i-1}\flowM \setminus F^i \flowM) \xrightarrow{\partial} \Omega_{k-1}(F^i \flowM) \to \dots
\]
Algebraically, these triangles fit together to form an \emph{exact couple}. The goal of this note is to:
\begin{enumerate}
    \item Define exact couples and derive the general algebraic formulas for their differentials.
    \item Translate these algebraic definitions into the concrete geometric language of flow modules and flow categories.
    \item Show how the higher differentials $d_{r+1}$ can be explicitly read off as the cobordism class of a closed manifold constructed via gluing.
\end{enumerate}

---

\section{Algebraic Theory of Exact Couples}

\begin{defn}
An \textbf{exact couple} consists of a pair of abelian groups (or objects in a stable $\infty$-category) $D$ and $E$, together with three homomorphisms:
\[
i \colon D \to D, \quad j \colon D \to E, \quad k \colon E \to D
\]
such that the diagram is exact at each vertex:
\[
\begin{tikzcd}
D \arrow[rr, "i"] & & D \arrow[dl, "j"] \\
& E \arrow[ul, "k"]
\end{tikzcd}
\]
Specifically:
\begin{enumerate}
    \item $\Image(i) = \operatorname{Ker}(j)$
    \item $\Image(j) = \operatorname{Ker}(k)$
    \item $\Image(k) = \operatorname{Ker}(i)$
\end{enumerate}
\end{defn}

An exact couple naturally defines a spectral sequence $\{E_r, d_r\}_{r \geq 1}$. 

\subsection{The First Page and Page 1 Differential}
The first page of the spectral sequence is defined by:
\[
E_1 \coloneqq E
\]
The differential on the first page $d_1 \colon E_1 \to E_1$ is defined as the composition:
\[
d_1 \coloneqq j \circ k
\]
\begin{prop}
The map $d_1$ satisfies $d_1^2 = 0$.
\end{prop}
\begin{proof}
By exactness at $D$, we have $\Image(j) = \operatorname{Ker}(k)$, which implies $k \circ j = 0$. Expanding $d_1^2$:
\[
d_1^2 = (j \circ k) \circ (j \circ k) = j \circ (k \circ j) \circ k = j \circ 0 \circ k = 0.
\]
\end{proof}

We can therefore define the second page $E_2$ as the homology of $E_1$ with respect to $d_1$:
\[
E_2 \coloneqq H(E_1, d_1) = \frac{\operatorname{Ker}(j \circ k)}{\Image(j \circ k)}
\]

\subsection{The $r$-th Page and Higher Differentials}
To define the higher pages and differentials, one algebraically constructs the \emph{derived couple} $(D_2, E_2, i_2, j_2, k_2)$ where:
\[
D_2 \coloneqq i(D) \subseteq D, \quad E_2 \coloneqq H(E, d_1)
\]
and the maps are given by $i_2 = i|_{D_2}$, $j_2(i(x)) = [j(x)]$, and $k_2([e]) = k(e)$. Iterating this construction yields the $r$-th page $E_r$ and its differential $d_r = j_r \circ k_r$. 

We can express these higher pages and differentials directly in terms of the original exact couple $(D, E, i, j, k)$ via the following theorem:

\begin{thm}\label{thm:algebraic_diff}
An element $e \in E_1$ survives to the $r$-th page $E_r$ if and only if:
\[
k(e) \in \Image(i^{r-1})
\]
Under this condition, let $y \in D$ be any element satisfying:
\[
k(e) = i^{r-1}(y)
\]
Then the differential of the class $[e] \in E_r$ is given by:
\[
d_r([e]) = [j(y)] \in E_r
\]
\end{thm}

\begin{proof}
We must prove that the formula $d_r([e]) = [j(y)]$ is well-defined:
\begin{enumerate}
    \item \textbf{Independence of the choice of $y$}: Suppose we have two elements $y, y' \in D$ such that $i^{r-1}(y) = k(e) = i^{r-1}(y')$. Then $i^{r-1}(y - y') = 0$, meaning $y - y' \in \operatorname{Ker}(i^{r-1})$. By induction and the exactness relations, any element in $\operatorname{Ker}(i^{r-1})$ is mapped by $j$ to the image of the boundaries from the previous pages. Thus:
    \[
    j(y) - j(y') = j(y - y') \in \Image(d_1) + \dots + \Image(d_{r-1})
    \]
    which means $[j(y)] = [j(y')]$ in $E_r$.
    \item \textbf{Compatibility with boundaries}: If $e$ is a boundary on an earlier page, say $e = d_1(e') = j(k(e'))$, then $k(e) = k(j(k(e'))) = 0$. We can therefore choose $y = 0$, which yields $d_r([e]) = [j(0)] = 0$.
\end{enumerate}
\end{proof}

---

\section{Geometric Case: Flow Categories and Proposition 2.20}

Let $\flowM$ be a flow category with a degree function $|\cdot| \colon \ob(\flowM) \to \bbZ$. Let $F^i \flowM \subset \flowM$ be the full subcategory consisting of objects of degree strictly greater than $i$:
\[
\ob(F^i \flowM) = \{x \in \ob(\flowM) \mid |x| > i\}
\]

\subsection{The Exact Couple of Bordism Groups}
For each $i \in \bbZ$, the subquotient category $F^{i-1}\flowM \setminus F^i \flowM$ consists of objects of degree exactly $i$. Since there are no non-trivial morphism spaces between objects of the same degree, this subquotient is a discrete category. Thus, a right flow module of dimension $k$ on $F^{i-1}\flowM \setminus F^i \flowM$ is simply a collection of closed, smooth, oriented manifolds $(\mathcal{N}(x))_{|x|=i}$ of dimension $k+i$.

From Hirschi's Proposition 2.20, we have exact triangles on the right bordism groups:
\[
\dots \to \Omega_k(F^i \flowM) \xrightarrow{\text{ext}} \Omega_k(F^{i-1} \flowM) \xrightarrow{\text{rest}} \Omega_k(F^{i-1}\flowM \setminus F^i \flowM) \xrightarrow{\partial} \Omega_{k-1}(F^i \flowM) \to \dots
\]
which we organize into the following exact couple:
\begin{itemize}
    \item $D = \bigoplus_{i, k} \Omega_k(F^i \flowM)$
    \item $E = \bigoplus_{i, k} \Omega_k(F^{i-1}\flowM \setminus F^i \flowM) \cong \bigoplus_{i, k} \bigoplus_{|x|=i} \Omega_{k+i}(pt)$
    \item $i \colon D \to D$ is the extension-by-$\emptyset$ map induced by the inclusion $F^i \flowM \hookrightarrow F^{i-1}\flowM$.
    \item $j \colon D \to E$ is the restriction-to-degree-$i$ map.
    \item $k \colon E \to D$ is the boundary map $\partial$ defined in Hirschi's Equation (2.3.1).
\end{itemize}

---

\subsection{The $d_1$ Differential}

Let $e = [\mathcal{N}] \in E_1$ be represented by a collection of closed manifolds $\mathcal{N}(x)$ of dimension $k+i$ for $|x|=i$. To find the differential $d_1([\mathcal{N}]) = j(k([\mathcal{N}]))$:
\begin{enumerate}
    \item \textbf{Apply $k$}: The boundary map $k([\mathcal{N}]) = \partial [\mathcal{N}]$ is the right flow module $\mathcal{N}'$ on $F^i \flowM$ (objects of degree $> i$) defined by:
    \[
    \mathcal{N}'(z) = \coprod_{|x|=i} \mathcal{N}(x) \times \flowM(x, z)
    \]
    for all $z \in \ob(\flowM)$ with $|z| > i$.
    \item \textbf{Apply $j$}: The restriction map $j$ projects this module to the next degree $i+1$. For any object $u$ of degree $i+1$, we obtain:
    \[
    (d_1([\mathcal{N}]))(u) = \coprod_{|x|=i} \mathcal{N}(x) \times \flowM(x, u)
    \]
\end{enumerate}
Since $\dim(\mathcal{N}(x)) = k+i$ and $\dim(\flowM(x, u)) = |u| - |x| - 1 = (i+1) - i - 1 = 0$, the spaces $\flowM(x, u)$ are $0$-dimensional manifolds (points with signs). The differential $d_1$ thus counts these isolated flow lines, recovering the standard Morse/Floer homology differential.

---

\subsection{Survival to Page $r+1$ and Flow Module Extensions}

By Theorem \ref{thm:algebraic_diff}, a class $[\mathcal{N}] \in E_{r+1}$ survives to page $r+1$ if and only if:
\[
k([\mathcal{N}]) \in \Image(i^r)
\]
Here, $i^r \colon \Omega_{k-1}(F^{i+r}\flowM) \to \Omega_{k-1}(F^i\flowM)$ is the inclusion map. Geometrically, this means that the boundary module $\mathcal{N}'$ on $F^i \flowM$ is equivalent in bordism to a module which is empty on objects of degree less than $i+r$.

This is equivalent to the statement that $\mathcal{N}'$ is null-bordant over the truncated subcategory containing objects of degree $i < |x| < i+r$. A null-bordism over this category is exactly a collection of manifolds with corners $\{\mathcal{W}(x)\}_{i < |x| < i+r}$ of dimension $k + |x|$ satisfying:
\[
\partial \mathcal{W}(x) = \left( \bigcup_{i < |z| < |x|} \mathcal{W}(z) \times \flowM(z, x) \right) \cup \left( \coprod_{|y|=i} \mathcal{N}(y) \times \flowM(y, x) \right)
\]
This boundary formula matches the exact definition of an \textbf{$r$-truncated right flow module} extending $\mathcal{N}$. We therefore obtain:

\begin{prop}[Survival Criterion]
A class $[\mathcal{N}] \in E_1$ survives to the $(r+1)$-th page of the spectral sequence if and only if $\mathcal{N}$ can be extended to an $r$-truncated right flow module $\mathcal{W}$ defined on all objects of degree up to $i+r-1$.
\end{prop}

---

\subsection{Reading Off the $d_{r+1}$ Differential}

Assuming $[\mathcal{N}]$ survives to the $(r+1)$-th page, we have our $r$-truncated right flow module $\mathcal{W}$ defined on objects of degree $i < |x| < i+r$. 

According to Theorem \ref{thm:algebraic_diff}, the differential $d_{r+1}([\mathcal{N}])$ is represented by $j(y)$ where $y$ is the extension of $\mathcal{W}$ to a module on $F^i \flowM$ (which is defined to be empty on objects of degree $\geq i+r$). The boundary of this extended module is a right flow module on $F^i \flowM$. 

Applying the restriction map $j$, we project this boundary to the next degree $i+r$. For any object $u$ of degree $i+r$, the differential is represented by the closed manifold:
\[
(d_{r+1}([\mathcal{N}]))(u) = \left( \bigcup_{i < |z| < i+r} \mathcal{W}(z) \times \flowM(z, u) \right) \cup \left( \coprod_{|y|=i} \mathcal{N}(y) \times \flowM(y, u) \right)
\]

\begin{thm}[Geometric Differential]
Let $[\mathcal{N}] \in E_{r+1}$. Let $\mathcal{W}$ be the right flow module extension of $\mathcal{N}$ on the degrees $i < |z| < i+r$. Then the differential $d_{r+1}([\mathcal{N}]) \in E_{r+1}$ is represented for each object $u$ of degree $i+r$ by the cobordism class of the closed manifold:
\begin{equation}\label{eq:geo_diff}
d_{r+1}([\mathcal{N}])(u) = \bigcup_{i < |z| < i+r} \mathcal{W}(z) \times \flowM(z, u) \cup \coprod_{|y|=i} \mathcal{N}(y) \times \flowM(y, u)
\end{equation}
of dimension $k+i+r-1$.
\end{thm}

This closed manifold in Equation \eqref{eq:geo_diff} is exactly the \textbf{Toda manifold} (or the geometric tensor product $W \circ \flowM_{-, i+r}$) formed by gluing the components of the extended right module $\mathcal{W}$ and $\mathcal{N}$ along the morphic moduli spaces of $\flowM$.

\end{document}
